\chapter{Future Work}
\label{cap:futuro}

The current evidence is limited to the specified zero-fee constant-product simulator. The following studies address the most important gaps before making claims about market operation.

\begin{enumerate}
  \item \textbf{Replay a documented market event.} Select one event with reproducible, redistribution-permitted data, such as the March 2023 USDC depeg, and align DEX swap, reserve and liquidity events with CEX bid/ask and depth records. Record block and transaction order, pool fees, gas, execution delay, token orientation, inventory constraints and data provenance. Freeze the event window, synchronization rule and metrics before analysis. Treat event replay as a case study, not broad validation; a second period or pool is needed to assess transfer.
  \item \textbf{Complete the liquidity-response factorial.} A shallow-depth pilot has now compared the reserve-linked arbitrage cap with a cap fixed to initial reserves; this substantially reduces the post-recovery E3 contrast but does not separate all mechanisms. Extend this control across depth, withdrawal rule, order scale and panic in a frozen factorial design. Include conditions where the withdrawal trigger does not fire, report recovery coverage as well as conditional post-recovery outcomes, and estimate seed-level interactions. Keep maximum deviation from the new reference, directional overshoot and post-recovery excess area as distinct outcomes.
  \item \textbf{Test an event-level flow-to-capacity diagnostic.} Replace configured order-size ratios with a prespecified measure based on realized signed flow, active reserves and available arbitrage capacity immediately before each event. Distinguish a single adverse impulse from sustained adverse flow. Vary flow direction, size, depth, capacity, withdrawal and panic; fit any candidate boundary on training cells and evaluate it on held-out depths, shocks and seeds against persistence and the raw capacity/depth inputs. The current ratios motivate this test but are not realized-flow measurements and do not establish a universal threshold.
  \item \textbf{Add execution costs and on-chain competition.} Implement pool fees, gas, settlement delay, finite trader capital and competition among arbitrage searchers. Relate these mechanisms to loss-versus-rebalancing and MEV models before interpreting capacity as an operational bottleneck.
  \item \textbf{Represent concentrated liquidity.} Extend the constant-product pool to positions bounded by price ranges, with active liquidity changing as price crosses those ranges. Keep the v2-style constant-product results as a deliberately narrow baseline.
  \item \textbf{Evaluate reduced models for a defined use.} Retain persistence and scalar AR(1) as baselines; compare them with a band-aware model using only information available at prediction time. Evaluate on chronological windows and separate pools. If the target is regime classification, define labels independently of the model inputs and measure classification performance directly.
  \item \textbf{Calibrate model time.} Give ticks an observable clock by linking events to timestamps or block intervals. Until that mapping is made, recovery is measured in event steps and cannot be interpreted as seconds, blocks or operational latency.
  \item \textbf{Evaluate engineering and pedagogy separately.} Test whether the interactive dashboard helps engineers or students distinguish a first return to the band from persistent tracking; this would assess the interface, not validate the market model.
\end{enumerate}
